\chapter{Securite, tests et deploiement}

\section{Introduction}
Ce chapitre presente les mecanismes de securite, la validation de la solution et son deploiement.

\section{Securite applicative}
Les mesures mises en place incluent :
\begin{itemize}
  \item hachage des mots de passe ;
  \item authentification JWT et verification de session ;
  \item validation des entrees (schema + sanitation) ;
  \item limitation du debit des requetes via Redis ;
  \item protection des routes administratives.
\end{itemize}

\section{Performance et scalabilite}
L'architecture exploite :
\begin{itemize}
  \item cache Redis pour reduire la latence ;
  \item appels non bloquants pour le tracking analytics ;
  \item streaming SSE pour ameliorer la reactivite percue ;
  \item approche modulaire favorisant la montee en charge.
\end{itemize}

\section{Strategie de tests}
La validation de la plateforme s'appuie sur :
\begin{itemize}
  \item tests des endpoints API ;
  \item tests fonctionnels par role ;
  \item tests de scenarios chat ;
  \item tests de non-regression sur les fonctionnalites critiques.
\end{itemize}

\section{Deploiement}
Le deploiement cible combine :
\begin{itemize}
  \item application Next.js en production ;
  \item PostgreSQL pour la persistence ;
  \item Redis pour cache et sessions ;
  \item variables d'environnement securisees.
\end{itemize}

\section{Conclusion}
Ces choix techniques garantissent une solution fiable, securisee et exploitable en contexte reel.
